% answers_short.tex —— 1977 年·福建卷·理科 · 简短答案速查
% 由 tools/make_short_answers.py 自 answers.tex 的 \gksolhead/\gkanswer 提取，
% 并与 1977/1977_数学真题_福建理_答案要点.md 逐题比对后生成，勿手改（改源 MD 后重跑）。

\begin{gknotice}
\textbf{使用说明}：本节只给最终答案便于快速核对，完整推导见解析版。答案为编写者独立推导，\textbf{非官方标准答案}。
\end{gknotice}

\gkpart{大题}
\begin{itemize}[leftmargin=2.6em]
  \item[\textbf{1.}] (1) $7$；(2) 负；(3) $(1,2)$；(4) $MN=2.5\ \mathrm{cm}$，$PQ=3\ \mathrm{cm}$；(5) $x=0.0003$；(6) $-1$；(7) $x=2$；(8) $\dfrac{a^n(a-3)}{a-1}$（在原式有意义的条件下）；(9) 最大值为 $\dfrac{49}{12}$，在 $x=-\dfrac56$ 时取得；(10) 见解析版；按源卷所示 V 形铁块作三视图草图
  \item[\textbf{2.}] (1) $-2<x<3$；(2) 见解析版；(3) $216\pi\ \mathrm g\approx678\ \mathrm g$（取 $\pi=3.14$）；(4) 在十一、十二月从十月产量按同一增长率递增的解释下，平均每月增长率为 $10\%$
  \item[\textbf{3.}] (1) 令 $q=\dfrac{\sqrt3}{2}$，则 $S_n=6R\dfrac{1-q^n}{1-q}$；(2) $S=\dfrac{6R}{1-q}=12R(2+\sqrt3)$
  \item[\textbf{4.}] $(x-5)^2+y^2=16$，为圆，圆心 $(5,0)$、半径 $4$。
  \item[\textbf{5.}] 按普通平面角定义，$AB=AC$ 且 $\angle BAC=60^\circ$ 与 $\angle ABP=120^\circ$、$\angle ACP=135^\circ$ 不相容，当前源面无法给出可靠唯一的 $AP$。
  \item[\textbf{6.}] 按源式可列 $\dfrac2\alpha-\tan\alpha=1$；若按弧度取锐角数值根，则 $\alpha\approx51.1361^\circ$，$m\approx11.0987$，$r\approx5.7786$，仅作按字面解释的候选。
  \item[\textbf{7.}] 源题面把数列 $1,2,4,\ldots$ 的前 $n$ 项和写成三次多项式，与 $S_n=2^n-1$ 不可能对所有 $n$ 同时成立；当前源面条件矛盾，待核。
\end{itemize}
